\documentclass[11pt,a4paper]{article}
\usepackage{acl_onecol}
\usepackage{amsmath}
\usepackage{array}
\usepackage{tabularx}
\usepackage{float}

\newcommand{\up}{$\uparrow$}
\newcommand{\down}{$\downarrow$}
\newcolumntype{L}{>{\raggedright\arraybackslash}X}

\title{RAG Assignment: Baseline Retrieval vs.\ Cross-Encoder Re-Ranking}

\author{Sankalp Chaturvedi \\
  2025201083}

\date{GitHub: \url{https://github.com/Sankalp0109/RAG_ASSIGNMENT_LMA}}

\begin{document}
\maketitle

\paragraph{Setup.} \emph{Document:} the 2024 FIA Formula One Sporting Regulations (Issue~6, 106 pages). It is a
regulatory text where most rules are followed by ``unless'', ``except'' or ``may not'' clauses, and the same
words (\emph{pit lane, penalty, tyres, safety car}) recur across many unrelated articles. \emph{Chunking:} one
chunk per sub-article (e.g.\ 55.8); long sub-articles are split at $\le$900 characters with a 150-character
overlap, giving 526 chunks. Both pipelines use this same chunk set. \emph{Bi-encoder:} \texttt{all-MiniLM-L6-v2}
with a FAISS inner-product index on normalised vectors (cosine similarity). \emph{Cross-encoder:}
\texttt{cross-encoder/}\allowbreak\texttt{ms-marco-MiniLM-L-6-v2}. \emph{LLM:} Gemini \texttt{gemini-3.6-flash}, temperature~0, with an
identical prompt for both pipelines. The prompt tells the model to answer only from the three passages, cite
articles, and say ``The provided context does not answer this'' otherwise. \emph{Pipeline~1} sends the bi-encoder
Top-3 of the Top-10; \emph{Pipeline~2} re-scores the same Top-10 with the cross-encoder and sends its Top-3.
\emph{Evaluation:} 52 questions, each labelled with the exact sentence in the PDF that answers it, and tagged as
keyword trap, exception, negation, buried chunk, or a direct-lookup control. All timings are on an Apple M4 (MPS).

\section{Architecture Comparison}

A \textbf{bi-encoder} encodes the query and each chunk \emph{independently} into one fixed vector
($\mathbf{q}=f(q)$, $\mathbf{d}=f(d)$) and scores them by cosine similarity $\cos(\mathbf{q},\mathbf{d})$. Since
$\mathbf{d}$ does not depend on the query, all chunks are embedded once, offline, and retrieval is just a
nearest-neighbour search. The limitation is that each chunk is compressed into a single 384-dimensional vector
before the question is known, and query and chunk only meet in one final dot product. No query token ever attends
to a passage token. Fine distinctions such as ``\emph{no} points are awarded'' vs.\ ``points are awarded'', or
``children'' vs.\ ``under the age of 16'', get averaged into a general topic vector.

A \textbf{cross-encoder} feeds the pair together, \texttt{[CLS] query [SEP] chunk [SEP]}, through the transformer
and predicts a relevance score from \texttt{[CLS]}. Self-attention runs \emph{jointly across the tokens of both
texts} in every layer, so each query word can attend to each passage word. The model can check whether the passage
actually answers the question (right condition, right polarity, right actor) instead of just sharing its topic. The
cost is that the score depends on the query, so nothing can be precomputed: every (query, chunk) pair needs a full
forward pass at query time.

\section{Qualitative \& Rank Shift Analysis}

\begin{table}[H]
\centering\small
\begin{tabular}{lrcccc}
\toprule
 & & \multicolumn{2}{c}{Hit@3 (gold reaches the LLM)} & \multicolumn{2}{c}{MRR@10} \\
\cmidrule(lr){3-4}\cmidrule(lr){5-6}
Query category & $n$ & Baseline & Re-ranked & Baseline & Re-ranked \\
\midrule
Negation / constraint      & 15 & 0.73 & \textbf{0.87} & 0.65 & \textbf{0.75} \\
Exception / corner case    & 13 & 0.85 & \textbf{0.92} & 0.64 & \textbf{0.83} \\
Keyword trap / distractor  & 10 & 0.60 & 0.60          & 0.50 & \textbf{0.57} \\
Buried chunk (rank 4--10)  &  3 & 0.33 & \textbf{0.67} & 0.20 & \textbf{0.73} \\
Control (direct lookup)    & 11 & 1.00 & 1.00          & \textbf{0.95} & 0.94 \\
\midrule
\textbf{All}               & 52 & 0.769 & \textbf{0.846} & 0.656 & \textbf{0.774} \\
\bottomrule
\end{tabular}
\caption{Rank of the gold chunk before and after re-ranking, per failure category. Hit@1 rises from 0.519 to 0.692.
The gold chunk is inside the baseline Top-10 for 47/52 queries.}
\label{tab:metrics}
\end{table}

Re-ranking leaves direct lookups unchanged, where the bi-encoder is already near-perfect. It helps most where the
answer depends on a condition or on polarity: negation (+14 points of Hit@3), exceptions (+7) and buried clauses.
Table~\ref{tab:shifts} lists every query whose gold chunk crossed the Top-3 boundary, i.e.\ every case where
re-ranking changed what the LLM was given. Table~\ref{tab:q10} shows one full Top-10 in detail, and
Table~\ref{tab:answers} quotes the two answers for three representative cases.

\begin{table}[H]
\centering\small
\begin{tabularx}{\textwidth}{@{}l>{\raggedright\arraybackslash}p{5.6cm}cL@{}}
\toprule
ID & Query (gold article) & Rank shift & Answer 1 (baseline) vs.\ Answer 2 (re-ranked) \\
\midrule
Q12 & Penalty for not using two dry tyre specs? (30.5m) & \#9 $\to$ \#1 \up & A1 \emph{misstates the rule}, gives an unrelated 54.3d penalty; A2 correct (DSQ / +30\,s) \\
Q10 & Must every driver use two different compounds? (30.5m) & \#8 $\to$ \#1 \up & A1 ``context does not answer''; A2 states the rule but omits the intermediates exception \\
Q50 & Penalty that cannot be served due to retirement? (54.3) & \#6 $\to$ \#1 \up & A1 refuses; A2 correct (grid penalty at next race) \\
Q26 & Can children enter the pit lane? (26.13) & \#5 $\to$ \#2 \up & A1 refuses; A2 gives the under-16 rule and time windows \\
Q21 & Can mechanics work on the car during a 5\,s penalty? (54.4c) & \#4 $\to$ \#2 \up & A1 right yes/no but cites the wrong clause (54.4d); A2 cites 54.4c \\
Q31 & When are no points awarded for a shortened race? (6.5a--b) & \#4 $\to$ \#2 \up & A1 refuses; A2 lists both no-points conditions \\
\midrule
Q03 & Top up fuel when the car pits for tyres? (36.2) & \#3 $\to$ \#5 \down & A1 \emph{correct}; A2 refuses (regression) \\
Q47 & How often must a rookie run in practice? (32.4c) & \#3 $\to$ \#5 \down & gold leaves the context (not sent to the LLM) \\
\bottomrule
\end{tabularx}
\caption{All queries whose gold chunk crossed the Top-3 boundary. Full Top-10 tables with bi-encoder and
cross-encoder scores are in the notebook (Section~7).}
\label{tab:shifts}
\end{table}

\begin{table}[H]
\centering\small
\begin{tabular}{@{}clcrcrl@{}}
\toprule
 & Chunk & Bi rank & Cosine & CE rank & CE score & Content \\
\midrule
 & 30.3 & 1 & 0.683 & 2  & 3.62  & Presenting tyres to the Technical Delegate \\
 & 30.5 & 2 & 0.680 & 5  & 2.19  & ``All tyres must be operated in accordance with \ldots'' \\
 & 30.5 & 3 & 0.653 & 4  & 2.26  & Same-specification sets; 54.3d penalty \\
 & 30.3 & 4 & 0.642 & 6  & 0.32  & Marking of the tyre sidewall \\
 & 30.5 & 5 & 0.607 & 7  & 0.16  & Additional tyres made available by the FIA \\
 & 30.6 & 6 & 0.582 & 10 & $-5.37$ & Testing of tyres \\
 & 30.2 & 7 & 0.581 & 9  & $-2.29$ & Tyre allocation per Competition \\
$\star$ & \textbf{30.5m} & \textbf{8} & \textbf{0.580} & \textbf{1} & \textbf{7.08} & \textbf{``must use at least two different specifications''} \\
 & 30.5 & 9 & 0.575 & 8  & $-1.29$ & Wet-weather tyres when a race is resumed \\
 & 30.2 & 10 & 0.571 & 3 & 2.59  & Allocation when extra sets are released \\
\bottomrule
\end{tabular}
\caption{Full Top-10 for Q10 (``Must every driver use two different tyre compounds in the race?''). Cosine scores
are almost flat (0.57--0.68), so the bi-encoder cannot tell the tyre chunks apart. The cross-encoder gives the gold
chunk ($\star$) a clear margin (7.08 vs.\ 3.62) and moves it \#8 $\to$ \#1.}
\label{tab:q10}
\end{table}

\paragraph{Exceptions \& corner cases (Q10, Q50).} The two-compound tyre rule sits inside the 5-page tyre article.
That article yields about ten near-identical chunks with cosine scores between 0.57 and 0.70, so the bi-encoder
orders them almost at random and the rule lands at \#8. The cross-encoder separates them sharply: the gold chunk
scores 7.1 against at most 3.6 for the rest, so it moves to \#1. Q50 is a true ``except when'' clause (a penalty
that cannot be served because the driver retired) hidden in the middle of the long 54.3 penalty list. It moves from
\#6 to \#1 with a cross-encoder score of $+5.5$, while every other candidate scores below $-0.2$.

\paragraph{Other: buried chunk (Q12).} This is the most harmful baseline case. The gold chunk sat at \#9, and the
baseline Top-3 contained a \emph{different} tyre penalty (54.3d, for not changing tyres as instructed). Answer~1
presented that penalty as the answer: a confident, cited, but wrong rule rather than a refusal. After re-ranking
(\#9 $\to$ \#1), Answer~2 gives the correct sanction: disqualification, or +30\,s if the race cannot be restarted.

\paragraph{Negation \& constraints (Q21, Q31).} For ``when are \emph{no} points awarded'', the bi-encoder's top hits
are the \emph{positive} points table (6.4) and the dead-heat rule (7.1). These are on-topic but have the wrong
polarity, and Answer~1 says the context does not answer. For Q21 (``may not be worked on'') the baseline Top-3 held
the neighbouring stop-and-go clause. Answer~1 reaches the right yes/no answer but cites the wrong clause, and adds
an engine-restart allowance that does not apply to a 5\,s penalty.

\paragraph{Keyword trap (Q26).} ``Children'' never appears in the document; the rule says ``no one under the age of
16''. The bi-encoder fills its Top-3 with chunks that merely share the words \emph{pit lane} (a car reversing in
the pit lane, the pit exit being closed). The cross-encoder connects the synonym and promotes the rule from \#5 to
\#2.

\begin{table}[H]
\centering\small
\begin{tabularx}{\textwidth}{@{}p{1.2cm}LL@{}}
\toprule
Query & Answer 1 (baseline Top-3) & Answer 2 (re-ranked Top-3) \\
\midrule
Q12 \newline \#9$\to$\#1 & ``According to Article 30.5, a penalty under Article 54.3d) will be imposed on any driver who does not
change tyres as specified.'' \newline \emph{Wrong rule: misstated penalty.} &
``Failure to use at least two different specifications of dry-weather tyres during a race will result in the
disqualification of the driver from the race results (Article 30.5). However, if the race is suspended and cannot be
re-started, \ldots\ thirty (30) seconds will instead be added \ldots'' \newline \emph{Correct.} \\
\midrule
Q31 \newline \#4$\to$\#2 & ``The provided context does not answer this.'' \newline \emph{Omission: gold not in context.} &
``\ldots\ no points are awarded if the leader has completed less than two laps (Article 6.5a). Additionally, no
points are awarded unless a minimum of two laps have been completed by the leader without a Safety Car and/or VSC
intervention (Article 6.5b).'' \newline \emph{Correct.} \\
\midrule
Q03 \newline \#3$\to$\#5 & ``According to Article 36.1, refuelling is only permitted in the Competitors' designated
garages. Furthermore, under Article 36.2, fuel may not be added to a car after it has left the pit lane \ldots\
until the end-of-session signal has been shown.'' \newline \emph{Correct.} &
``The provided context does not answer this.'' \newline \emph{Regression: re-ranking removed the gold chunk.} \\
\bottomrule
\end{tabularx}
\caption{Verbatim LLM answers (identical prompt, only the Top-3 chunks differ; ``\ldots'' marks trimmed text). All 14
answers are in the notebook.}
\label{tab:answers}
\end{table}

\paragraph{Where re-ranking did not help.} (i) In Q03 the question mixes two topics (fuel \emph{and} tyres). The
cross-encoder followed \emph{tyres} and pushed the refuelling chunk out of the Top-3, turning a correct Answer~1
into a refusal. All of its scores for this query were strongly negative ($-5.6$ to $-10.7$), so it was not
confident about any candidate. (ii) For 5 queries the gold chunk was not in the Top-10 at all, so a re-ranker over
the Top-10 cannot recover it. (iii) In Q10 the chunk boundary splits clause 30.5m, so the words ``Unless he has
used intermediate or wet-weather tyres'' are in the previous chunk. Answer~2 therefore states the rule without its
exception: re-ranking fixes ordering, not chunking.

\section{Latency Profiling}

Retrieval stages were timed over all 52 queries $\times$ 20 repeats (1{,}040 runs, after warm-up). LLM time is the
wall-clock time of the 14 Gemini calls in the failure-query analysis.

\begin{table}[H]
\centering\small
\begin{tabular}{lrrr}
\toprule
Stage & Mean (ms) & p50 (ms) & p95 (ms) \\
\midrule
Query embedding (bi-encoder)                 & 5.32  & 5.30  & 5.66 \\
FAISS Top-10 search (526 vectors)            & 0.02  & 0.02  & 0.03 \\
Cross-encoder re-rank (10 pairs)             & 44.70 & 45.07 & 50.13 \\
\midrule
\textbf{Pipeline 1 retrieval (baseline)}     & \textbf{5.34}  & 5.32  & 5.69 \\
\textbf{Pipeline 2 retrieval (re-ranking)}   & \textbf{50.04} & 50.47 & 55.55 \\
\midrule
LLM generation (Gemini API, per call)        & 6{,}526 & 5{,}275 & 12{,}231 \\
Pipeline 1 end-to-end                        & 6{,}531 & 5{,}280 & 12{,}236 \\
Pipeline 2 end-to-end                        & 6{,}576 & 5{,}325 & 12{,}281 \\
\bottomrule
\end{tabular}
\caption{Latency per query (Apple M4, MPS). Embedding all 526 chunks and building the index is a one-off 2.0\,s.}
\label{tab:latency}
\end{table}

Re-ranking makes retrieval about $9\times$ slower (5\,ms $\to$ 50\,ms), but in absolute terms it adds only 45\,ms.
That is about 0.7\% of end-to-end latency, because the LLM call (5--12\,s) dominates everything. Ten
cross-encoder passes (about 4.5\,ms per pair) are a small price for the gains in Table~\ref{tab:metrics}.

\section{Scaling Question}

\begin{table}[H]
\centering\small
\begin{tabular}{rccc}
\toprule
Chunks & Bi-encoder (flat search) & Two-stage (+ re-rank Top-10) & Cross-encoder on \emph{all} chunks \\
\midrule
526 (this corpus) & 5.3 ms  & 49.6 ms & 3.3 s \\
10,000            & 5.7 ms  & 50.0 ms & 1.1 min \\
100,000           & 9.7 ms  & 53.9 ms & 10.5 min \\
1,000,000         & 49.1 ms & 93.4 ms & 1.8 h \\
10,000,000        & 444 ms  & 488 ms  & 17.6 h \\
\bottomrule
\end{tabular}
\caption{Per-query cost. The first row is measured; larger rows extrapolate linearly from the measured costs
(6.3\,ms per cross-encoder pair, and flat-search time scaled with corpus size).}
\label{tab:scaling}
\end{table}

We do not run the cross-encoder over the whole collection because its cost grows linearly with the collection and
cannot be moved offline. To test this, I scored all 526 chunks with the cross-encoder for a single query: it took
3.3\,s, about 600$\times$ the bi-encoder retrieval. A cross-encoder score exists only for a specific
(query, chunk) pair, so for \emph{every} new question the full transformer must run once per chunk. With a
bi-encoder, the chunk-side work (embedding) is done once, ahead of time. Each query then costs one short encoding
plus a vector search, and an approximate-nearest-neighbour index can keep that search fast even for millions of
vectors. Extrapolating, a cross-encoder over one million chunks would take about 1.8\,h per question, compared with
under 0.1\,s for the two-stage pipeline (Table~\ref{tab:scaling}).

Scoring everything would also add little accuracy, because the cheap first stage already has high recall (the gold
chunk was in the Top-10 for 47 of 52 queries). The sensible design is therefore a funnel. The fast, high-recall
bi-encoder narrows the collection to a few candidates, and the slow, high-precision cross-encoder is spent only on
ordering those candidates. That is exactly where it changed the answers in Table~\ref{tab:shifts}.

\end{document}
